\chapter{绪论}
\section{研究背景}
移动边缘计算通过在接入侧配置算力，减少终端任务到云中心的往返时延，但无线链路和计算资源之间存在紧密耦合~\cite{mao2017survey}。无人机具备按需部署和视距传播优势，可在灾害救援、临时热点和基础设施薄弱区域形成空中接入与边缘计算节点~\cite{zeng2016uav,alhourani2014optimal}。当多架无人机复用全部频谱时，系统获得更高的频谱利用率，同时也产生跨无人机同频干扰。

上行分组 NOMA 允许两个用户复用一个子载波，并由接收端执行串行干扰消除~\cite{ding2017survey}。用户关联、子载波位置和强弱角色共同决定接收信干噪比。若一个用户从无人机 $k$ 调整到 $k'$，其服务信道发生变化，同时会改变其他无人机在相同子载波上的干扰。因此，独立优化每架无人机可能遗漏网络级收益。

另一方面，终端任务具有数据量、计算周期和截止时间。只最大化通信速率并不能保证按时完成；在有限 MEC 算力下，控制器还需确定最小卸载比例、最小边缘 CPU 需求及任务准入顺序。

\section{问题边界与专业硕士价值}
本文研究固定高度、固定悬停位置、集中式控制的时隙调度。优化变量包括任务的实际服务无人机、子载波、NOMA 角色、卸载比例和 MEC CPU。无人机轨迹、发射功率、电池、长期队列和现场部署不属于本文优化变量。

专业硕士价值体现在完整工程链路而非单一公式：显式调度数据结构、真实服务接收机速率计算、硬 SIC 检查、同任务公平基线、可失败的初始化、确定性搜索、原子实验产物、配对统计检验以及一键复现。每个到达任务只计数一次，方法间共享不可变外生轨迹。

\section{两项创新}
本文只定义两项创新：
\begin{enumerate}
 \item 干扰外部性感知的跨 UAV 联合关联与配对。通过 Replace、Position 和 CrossUAV 邻域改变调度拓扑，并用全网词典序目标接受严格改进；每个候选均重新计算实际服务信道、跨无人机干扰、SIC 与 CPU。
 \item 面向截止时间的卸载与 MEC 算力分配。由截止时间推导每项任务的最小卸载比例与最小 MEC CPU，按需求升序选择最长可行前缀，从而最大化单架无人机按时完成任务数。
\end{enumerate}
二者分别改变“调度拓扑/接受准则”和“卸载/算力准入”，可独立消融后组合。

\section{论文结构}
第二章综述相关工作；第三章建立唯一任务语义、物理层和优化目标；第四、第五章分别给出两项创新；第六章说明实现、实验协议和正式结果；第七章讨论部署与局限；第八章总结全文。
